\begin{tikzpicture}

	%%% in silico evolution scheme %%%
	\coordinate[yshift = -\columnsep] (evo scheme) at (0, 0);
	
	\node[anchor = north west, xshift = .5\columnsep] (spacer) at (evo scheme) {\hphantom{computational model}};
	
	\node[anchor = north, align = center, font = \bfseries, yshift = -.5\columnsep] (evo title) at (spacer.north) {\textit{in silico} evolution};
	
	\node[below = .1cm of evo title, align = center, inner ysep = .2em] (start) {WT sequence};
	
	\node[below = .5cm of start, align = center, inner ysep = .2em] (mutate) {digitally create all\\possible single-base\\substitution variants};
	
	\node[below = .5cm of mutate, align = center, inner ysep = .2em] (score) {score variants with\\\textbf{plantGREP}};
	
	\node[below = .5cm of score, align = center, inner ysep = .2em] (select) {select sequence\\with best score};
	
	\draw[thick, -Latex] (start) -- (mutate);
	
	\draw[thick, -Latex] (mutate) -- (score);
	
	\draw[thick, -Latex] (score) -- (select);
	
	\draw[thick, -Latex, rounded corners = .15cm] (select.east) -- ($(select -| spacer.east) + (\columnsep, 0)$) |- (mutate.east) node[anchor = south, rotate = -90, pos = .25, align = center] (n rounds) {$n$ rounds};
	
	
	%%% in silico evolution: number of rounds %%%
	\coordinate[xshift = \columnsep] (evo rounds) at (evo scheme -| n rounds.north);
	
	\distance{evo rounds}{\textwidth, 0}
	
	\begin{pggfplot}[%
		at = {(evo rounds)},
		total width = \xdistance,
		xlabel = number of \textit{in silico} evolution rounds,
		ylabel = \enrichment,
		enlarge x limits = 0,
		enlarge y limits = {lower = 0.1, upper = 0.1},
		title is csname,
		title color from name,
		xytick = {-12, -10, ..., 12},
	]{evolution_rounds}
	
		\pggf_annotate(zeroline)
	
		\pggf_violin(
			color from col name,
			draw = black,
			shade = left,
			shade color = white,
			cld = {anchor = south, text depth = 0pt}{max}
		)
	
	\end{pggfplot}
	
	
	%%% in silico evolution: species/condition specificity %%%
	\coordinate[yshift = -\columnsep] (evo specificity) at (evo scheme |- xlabel.south);
	
	\colorlet{start}{white}

	\tikzset{
		species-specificity/.style = {left color = tobacco!20, right color = maize!20, middle color = white},
		light-specificity/.style = {left color = light!20, right color = dark!20, middle color = white},
		temperature-specificity/.style = {left color = warm!20, right color = cold!20, middle color = white},
		TRUE/.style = {fill = pgffillcolor!50!white},
		FALSE/.style = {}
	}
	
	\expandafter\def\csname species-specificity\endcsname{species specificity}
	\expandafter\def\csname light-specificity\endcsname{light specificity}
	\expandafter\def\csname temperature-specificity\endcsname{temperature specificity}
	
	\begin{pggfplot}[%
		at = {(evo specificity)},
		total width = \textwidth,
		xlabel = number of \textit{in silico} evolution rounds,
		ylabel = $\log_2$(specificity),
		enlarge x limits = {abs = 1, rel = 0.005},
		enlarge y limits = {lower = 0.1, upper = 0.1},
		title is csname,
		title style from name,
		ytick = {-12, -10, ..., 12},
		xticklabels = {12, 6, 0, 6, 12},
		legend columns = 2,
		legend style = {anchor = north east, at = {(1, 1)}},
		legend image width = .75\baselineskip,
		tight y
	]{evolution_specificity}
	
		\pggf_annotate(
			zeroline,
			manual legend = {\textbf{objective:} = {empty legend}},
		)
		
		\pggf_annotate(
			annotation = {
				\draw[thick, -Latex] (-.25, .25) -- (-.25, 8.5) node[pos = .45, anchor = north, rotate = 90, align = center, font = \figsmall] {more active\\in \textbf{\usename{tobacco}}};
				\draw[thick, -Latex] (6.25, -.25) -- (6.25, -8.5) node[pos = .45, anchor = north, rotate = -90, align = center, font = \figsmall] {more active\\in \textbf{\usename{maize}}};
			},
			manual legend = {
				tobacco-specific = {violin legend = y, fill = tobacco, fill opacity = .5},
				\relax = {empty legend},
				maize-specific = {violin legend = y, fill = maize, fill opacity = .5}
			},
			only facet = 1
		)
		
		\pggf_annotate(
			annotation = {
				\draw[thick, -Latex] (-.25, .5) -- (-.25, 8.5) node[pos = .45, anchor = north, rotate = 90, align = center, font = \figsmall] {more active\\in \textbf{\usename{light}}};
				\draw[thick, -Latex] (6.25, -.5) -- (6.25, -8.5) node[pos = .45, anchor = north, rotate = -90, align = center, font = \figsmall] {more active\\in \textbf{\usename{dark}}};
			},
			manual legend = {
				light-specific = {violin legend = y, fill = light, fill opacity = .5},
				\relax = {empty legend},
				dark-specific = {violin legend = y, fill = dark, fill opacity = .5}
			},
			only facet = 2
		)
		
		\pggf_annotate(
			annotation = {
				\draw[thick, -Latex] (-.25, .5) -- (-.25, 8.5) node[pos = .45, anchor = north, rotate = 90, align = center, font = \figsmall] {more active\\in \textbf{\usename{warm}}};
				\draw[thick, -Latex] (6.25, -.5) -- (6.25, -8.5) node[pos = .45, anchor = north, rotate = -90, align = center, font = \figsmall] {more active\\in \textbf{\usename{cold}}};
			},
			manual legend = {
				warm-specific = {violin legend = y, fill = warm, fill opacity = .5},
				\relax = {empty legend},
				cold-specific = {violin legend = y, fill = cold, fill opacity = .5}
			},
			only facet = 3
		)
	
		\pggf_violin(
			style from column = {fill = color, fade},
			cld = {anchor = south, text depth = 0pt}{max},
			forget plot
		)
	
	\end{pggfplot}
	
	
	%%% subfigure labels %%%
	\subfiglabel{evo scheme}
	\subfiglabel{evo rounds}
	\subfiglabel{evo specificity}

\end{tikzpicture}